% Includable section; requires amsmath, amssymb. Labels have prefix B:.
\section{Spreading the source forms and choosing the determinant family}
\label{B:section}

In the notation of \S\ref{sz:setup}, put \(Q=\rho(e_0)\)
for the good-type projector, so that \(Y_T=QY\).
On the full source the maps associated to the residue action are
\[
  G=\prod_{i=0}^1\prod_{w\mid p_i\text{ chosen}}\mathrm{GL}_r(\mathbb F_w),
  \qquad B_g=|G|Q\rho(g)\Pi:X\longrightarrow Y_T.
\]
The restriction of \(B_g\) to \(S\), multiplied by the nonzero
scalar \(c_B c_{\rm ss}^{-1}\) of Theorem~\ref{A:main}, is \(P_g\).

\paragraph{Theorem B (spreading).}
Under hypotheses F0--F7 of \S\ref{B:frontier},
\begin{equation}
\label{B:theorem}
 d_T>0,\qquad
 |t|d_S\le5d_T\log(1+2b^2),\qquad
 \dim\sum_{g\in t}P_g(S)\ge(1-4/b^2)d_T
\end{equation}
for one finite set \(t\subseteq G\).
The arithmetic hypotheses give the two estimates
\eqref{B:near} and \eqref{B:moment} on the same scaled map.

\subsection{From two Hilbert--Schmidt estimates to a new translate}
\label{B:one}
Choose the source coordinates indexed by the pairs consisting of a
source class and a Bessel coordinate (the coordinate isomorphism
includes the central translation by the lattice scale).
In these coordinates \(X\) is a nonzero Euclidean space. Define
\({\Pi_{\rm sc}}=(\operatorname{vol}(K_{\rm pr})\nu)^{-1/2}\Pi\) in these coordinates,
where \(K_{\rm pr}\) is the principal level and \(\nu>0\) is the unipotent
covolume. Both volume and covolume are positive. We use the Hilbert--Schmidt norm
\(\|C\|_{\rm HS}^2=\operatorname{tr}(C^*C)\).
The estimates furnished by the arithmetic argument are
\begin{align}
 \|Q{\Pi_{\rm sc}}\,\operatorname{pr}_Z\|_{\rm HS}^2
 &\ge \dim Z-5b^{-4}d_X\quad(Z\subset X), \label{B:near}\\
 d_T\sum_{g\in G}\|{\Pi_{\rm sc}}^*\rho(g){\Pi_{\rm sc}}\|_{\rm HS}^2
 &\le |G|(1+9b^{-4})d_X^2. \label{B:moment}
\end{align}
Their reduction to more primitive inputs is given below.
For \(Z=X\), \eqref{B:near} gives \(\|Q{\Pi_{\rm sc}}\|_{\rm HS}^2>0\),
since \(b\ge1000\) and \(d_X>0\). Hence \(d_T>0\).
Set \(\epsilon=(8b^2)^{-1}\). The admissible scales imply
\(5b^{-4},9b^{-4}\le\epsilon\) and \(0<\epsilon<1/16\).

Let \(F\subset Y_T\), \(\alpha=\dim F/d_T\), and
\(C_g=\rho(g){\Pi_{\rm sc}}{\Pi_{\rm sc}}^*\rho(g)^*\). By cyclicity of trace,
\[
 \langle C_g,C_h\rangle_{\rm HS}
 =\|{\Pi_{\rm sc}}^*\rho(g^{-1}h){\Pi_{\rm sc}}\|_{\rm HS}^2 .
\]
Summing over both indices and tracing against \(\operatorname{pr}_F\)
gives the identities
\[
 \left\|\sum_g C_g\right\|_{\rm HS}^2
 = |G|\sum_g\|\Pi_{\rm sc}^*\rho(g)\Pi_{\rm sc}\|_{\rm HS}^2,
 \qquad
 \operatorname{tr}(\operatorname{pr}_F C_g)
 = \|\operatorname{pr}_F\rho(g)\Pi_{\rm sc}\|_{\rm HS}^2.
\]
Cauchy--Schwarz in the space of operators and \eqref{B:moment} give
\[
 \left(\frac1{|G|}\sum_g
     \|\operatorname{pr}_F\rho(g){\Pi_{\rm sc}}\|_{\rm HS}^2\right)^2
 \le\alpha(1+\epsilon)d_X^2 .
\]
Since \(2\sqrt{\alpha(1+\epsilon)}\le\alpha+1+\epsilon\),
some \(g\) satisfies
\(\|\operatorname{pr}_F\rho(g){\Pi_{\rm sc}}\|_{\rm HS}^2
\le(1+\alpha+\epsilon)d_X/2\).
Put \(Z=(Q\rho(g){\Pi_{\rm sc}})^{-1}F\).
On \(Z\) the maps \(Q\rho(g){\Pi_{\rm sc}}\) and \(\operatorname{pr}_F\rho(g){\Pi_{\rm sc}}\)
agree: the component orthogonal to \(Y_T\) is orthogonal to \(F\).
Moreover \(Q\) commutes with the unitary \(\rho(g)\).
Thus
\[
 \dim Z-\epsilon d_X
 \le\|Q{\Pi_{\rm sc}}\operatorname{pr}_Z\|_{\rm HS}^2
 =\|\operatorname{pr}_F\rho(g){\Pi_{\rm sc}}\operatorname{pr}_Z\|_{\rm HS}^2
 \le(1+\alpha+\epsilon)d_X/2.
\]
Coordinate change and multiplication of a map by a nonzero scalar
preserve the dimension of the preimage of every subspace. In particular
\(Q\rho(g){\Pi_{\rm sc}}\) and \(B_g\) have preimages of equal dimension.
We have proved, for every \(F\subset Y_T\), the existence of \(g\) with
\begin{equation}
\label{B:spread}
 \dim B_g^{-1}F\le\left(\frac{1+\alpha}{2}
                  +\frac32\epsilon\right)d_X .
\end{equation}
The argument uses the full unprojected moment
\({\Pi_{\rm sc}}^*\rho(g){\Pi_{\rm sc}}\) and the projected inferior estimate for \(Q{\Pi_{\rm sc}}\).

\subsection{Greedy selection with a cardinality bound}
\label{B:greedy}
If \(c=d_T-\dim F>16\epsilon d_T\), apply \eqref{B:spread}
to \(F\). The map to \(Y_T/F\) induced by \(B_g\) has kernel
\(B_g^{-1}F\); therefore
\[
 \dim(F+B_g(X))-\dim F
 \ge (\tfrac12(1-\alpha)-\tfrac32\epsilon)d_X
 \ge \tfrac25(1-\alpha)d_X.
\]
Indeed \(\tfrac12-\tfrac3{32}=\tfrac{13}{32}>\tfrac25\).
Write \(\lambda=2d_X/(5d_T)>0\). Starting with \(F=0\), the first
step is possible, and its gain cannot exceed \(d_T\), so \(\lambda\le1\).
At each step where the codimension exceeds \(16\epsilon d_T\), a
translate reduces it to at most \((1-\lambda)c\le e^{-\lambda}c\).
Such a translate cannot have been selected previously (the gain is
positive). Stop if the threshold is reached, otherwise use at most
\[
 j=\left\lceil\lambda^{-1}\log(1/(16\epsilon))\right\rceil
\]
steps. After \(j\) steps the missing share is at most
\(e^{-j\lambda}\le16\epsilon\); stopping earlier also suffices.
Since \(j\lambda\le1+\log(1/(16\epsilon))\), the set \(t\)
so obtained satisfies simultaneously
\begin{align}
 |t|d_X&\le\tfrac52(1+\log(1/(16\epsilon)))d_T
 \le5d_T\log(1+2b^2),\label{B:count}\\
 \dim\sum_{g\in t} B_g(X)&\ge(1-16\epsilon)d_T
 =(1-2/b^2)d_T.\label{B:full}
\end{align}
For the second inequality of \eqref{B:count}, use
\(1/(16\epsilon)=b^2/2\le1+2b^2\) and
\(\log(1+2b^2)\ge1\). The argument works for the stated dimension
ratios; no implicit rounding of the number of translates is used.

\subsection{The loss from restricting to ramified forms}
\label{B:ramification}
The scalar translations at the \(f\) places dividing \(\pi\) have
order dividing \(N-1\). By the free-level dimension count
in \eqref{A:dimension}, the intersection defining \(S\) has
\begin{equation}
\label{B:codim}
 d_X-d_S\le \frac{f}{N-1}d_X\le b^{-8}d_X,
\end{equation}
using \(f\le b\) and \(N-1\ge b^9\).

Choose a vector-space complement \(X=S\oplus C\). For the
\emph{same} set \(t\), the dimension of
\(\sum_g B_g(X)\) is at most that of \(\sum_g B_g(S)\) plus
\(|t|\dim C\). From \eqref{B:count}, \eqref{B:codim},
\[
 |t|(d_X-d_S)\le 5 b^{-8}\log(1+2b^2)d_T
 \le (2/b^2)d_T .
\]
The last inequality holds, for example, by
\(\log(1+2b^2)\le 2b^2\) and \(b^4\ge5\).
Multiplying the restricted maps by \(c_Bc_{\rm ss}^{-1}\) does not
change their ranges. Therefore
\[
 \sum_{g\in t}P_g(S)=\sum_{g\in t}B_g(S),\qquad
 \dim\sum_{g\in t}P_g(S)\ge(1-4/b^2)d_T.
\]
As \(d_S\le d_X\), \eqref{B:count} also bounds \(|t|d_S\).
This proves \eqref{B:theorem} from \eqref{B:near}--\eqref{B:moment}.

\subsection{Arithmetic origin of the two moment estimates}
\label{B:arithmetic}
We now trace \eqref{B:near} and \eqref{B:moment} further.
Use the filling data of \(D\), write \(M=p_0p_1\), and let
\(R\subset G=G_0\times G_1\) be the reduction of the Bessel subgroup.
Let \(\delta\) be the product of the local Bessel formal degrees and
\(\mu\) the mass of the genus of the standard definite rank-\(r\)
lattice. Put
\[
 I_0=|R|d_X,\qquad I_*=d_X^2/(\delta\mu),\qquad q_g={\Pi_{\rm sc}}^*\rho(g){\Pi_{\rm sc}}.
\]
The class representatives \(h_c\), cutoffs \(\beta_c\) and vector
Bessel function \(W\) are those in the filling data. Write
\(\mathcal A(g;c,c')\) for the matrix
\[
 \sum_{\gamma:\ h_c^{-1}\gamma h_{c'}\in K,\,
             \operatorname{red}(h_c^{-1}\gamma h_{c'})=g}
 \int \beta_c(x)\beta_{c'}(\gamma^{-1}x)
             W(x)W(\gamma^{-1}x)^*\,dx .
\]
Here \(K\) is the lattice stabilizer; the integral is on the
product of real unitary groups in the filling model. The sum is
over rational points. In the source orthonormal basis indexed
by \((c,i)\), the blocks of \(q_g\) are these blocks divided by
\(\nu\), with entrywise conjugation if the inner product is taken
linear in its second variable. In particular
\(\|q_g\|_{\rm HS}^2=\nu^{-2}\sum_{c,c'}\|\mathcal A(g;c,c')\|^2\).

The unfolding giving this identity begins with the Petersson
product of two Poincar\'e series, unfolds the series belonging to
\((c,i)\) to its seed, and evaluates the other series at the
class point \(h_c\) translated by a representative of \(g\).
The finite support of the seed selects exactly the displayed
congruence condition, with factor \(\operatorname{vol}(K_{\rm pr})\).
The real-homeomorphism change of variables produces the product
of the two cutoffs and Bessel functions. Scaling both series
by \((\operatorname{vol}(K_{\rm pr})\nu)^{-1/2}\) gives the factor
\(\nu^{-1}\). F1 is the fitted identification of the period-defined \(\Pi\)
with this series and the normalized seed and block identities.
The absolute summability of the amplitude integrals used here
follows from the majorization in \S\ref{B:cutoff}.

\paragraph{Near the Bessel subgroup.}
Denote by \(r_\psi(g)\) the operator whose blocks come from the
same amplitude sum restricted to the Bessel subgroup; it is a
unitary representation on \(R\), extended by zero off \(R\).
Here is how the action and the trace enter the proof.
For \(t\in R\), a rational Bessel term over the column class and
\(t^{-1}\) specifies a new class and a unitary change of Bessel
coordinates. In the Bessel-only sum on that column, the block
at the new class is this unitary matrix (with the conjugation
fixed by the inner product), and every other block vanishes.
Indeed one moves the term to the identity block, where integrating
the two cutoffs and the Bessel function gives \(\nu\) times the
identity. This identity-block normalization and the change of class
coordinates are among the fitted identities in F1. Composing two
term changes multiplies their matrices. The blocks at inverse group
elements are adjoints by replacing the rational term by its inverse.
Thus \(r_\psi(1)=1\) and \(r_\psi\) is unitary on \(R\).
Applying the same change of variables on the right to the
full sum gives \(q_{gt}=q_g r_\psi(t)\). For \(u\in X\),
\(d=\rho(t){\Pi_{\rm sc}}u-{\Pi_{\rm sc}}r_\psi(t)u\) now satisfies
\({\Pi_{\rm sc}}^*\rho(g)^{-1}d=0\) for every \(g\); hence it is orthogonal to
the span of every \(\rho(g){\Pi_{\rm sc}}X\). It belongs to that very span, so
\(d=0\). This proves \(\rho(t){\Pi_{\rm sc}}={\Pi_{\rm sc}}r_\psi(t)\).
Let \(E_g=q_g-r_\psi(g)\); thus \(E_{sg}=r_\psi(s)E_g\) and
\(E_{gs}=E_gr_\psi(s)\) for \(s\in R\). Write
\(\Omega_i=R G_i\), where \(G_i\) is embedded as one factor.
The quantitative estimate needed is
\begin{equation}
\label{B:nearBessel}
 \sum_{g\in\Omega_i}\|E_g\|_{\rm HS}^2\le b^{-12}I_0
 \qquad(i=0,1).
\end{equation}
For \(g\in R\), \(E_g=({\Pi_{\rm sc}}^*{\Pi_{\rm sc}}-1)r_\psi(g)\), so its squared
Hilbert--Schmidt norm is constant. Restricting \eqref{B:nearBessel}
to \(R\) and dividing by \(|R|>0\) yields
\begin{equation}
\label{B:defect}
 \|{\Pi_{\rm sc}}^*{\Pi_{\rm sc}}-1\|_{\rm HS}^2\le b^{-12}d_X .
\end{equation}
For every subspace \(Z\subset X\),
\[
 \|{\Pi_{\rm sc}}\operatorname{pr}_Z\|_{\rm HS}^2
 =\dim Z+\operatorname{tr}(\operatorname{pr}_Z({\Pi_{\rm sc}}^*{\Pi_{\rm sc}}-1))
 \ge\dim Z-b^{-6}d_X\ge\dim Z-b^{-4}d_X.
\]
Also \(\|q_1\|_{\rm HS}^2\le(1+b^{-6})^2 d_X\le 2d_X\).

To prove \eqref{B:nearBessel}, group the terms by links.
For a subset of reductions \(\Omega\),
let \(\mathcal E(\Omega)\) be the sum of squared cut-amplitude norms, taken
link by link for terms of nonzero frame discriminant. A link
is the submodule of pairs of classes of \(v,\gamma^{-1}v\) in
the quotient of the orthogonal complement of the isotropic frame
span by that span, with \(v\) orthogonal to the frame span and
its \(\gamma\)-translate. Let \(\operatorname{meetNorm}(g)\)
be the product of the norms of places where the frame span
meets its translate nontrivially. The near-pair count gives the following bound. If \(0<z\le\operatorname{meetNorm}(g)\) throughout \(\Omega\),
\[
 \mathcal E(\Omega)/\nu^2\le (b^{2f}/z)^k I_*;
\]
the error from different link classes/zero discriminant is
\(b^{-kf}I_0\) in the normalized cut second moment and the
error of removing the cut is at most \(b^{-40}I_0\) in squared
norm (F2--F4). On a region avoiding nontrivial root-of-unity
scalar cosets the off-Bessel terms equal the remaining terms.
Expanding squares by links, then taking the triangle inequality
in the space of all blocks, gives
\begin{equation}
\label{B:rootbound}
 \sqrt{\sum_{g\in\Omega}\|E_g\|_{\rm HS}^2}
 \le\sqrt{\mathcal E(\Omega)/\nu^2}+ b^{-15}\sqrt{I_0},
\end{equation}
since \(kf\ge40\), and \(2b^{-20}\le b^{-15}\).
For \(\Omega_i\) the reductions match a Bessel element at the
other prime and its conjugate; the scalar-coset avoidance
follows from the first frame vector, and
\(z\ge(M/b)^f\) by the comparability of the level primes.
The filling bounds \(I_*\le(M/b^6)^{kf}I_0\) and \(kf\ge40\)
give
\[
 (b^{2f}/z)^k I_*\le b^{-3kf}I_0\le b^{-30}I_0.
\]
Thus the right side of \eqref{B:rootbound} is at most
\(2b^{-15}\sqrt{I_0}\); squaring and using \(4b^{-30}\le b^{-12}\)
proves \eqref{B:nearBessel}.

\paragraph{The cost of projecting to good types.}
At prime \(p_i=q\), each of the \(f\) fields has \(q\) elements.
The prime-wise good idempotent \(e_i^{\rm pr}\) is the sum over tuples of labels
of degree \(r\) that are allowed at every place and primitive at
at least one place. Allowed means monic squarefree with no root
at zero or at an \(M_0\)-th root of unity in the field, where
\(M_0=q_2\) is the exclusion modulus of \S\ref{sz:setup}.
Primitive means irreducible with a root of order \(q^r-1\).
For each irreducible factor \(p\) of a squarefree label, write
\(d=\deg p\) and choose a root in \(\mathbb F_{q^d}^\times\).
The faithful identification of this cyclic group with its dual
commutes with raising to \(q\), so the degree-\(d\) Frobenius
orbit of roots gives a degree-\(d\) orbit of characters. The
Hall function \(\tau_p\) used in the level projector equals the one-box Green
character of this orbit, in the positive-degree normalization, by
Lemma~\ref{lem:fitted-primary-character}.
Green's theorem~\cite[\S8, Theorem~14 and p.~444]{green}
gives its degree \(\prod_{h=1}^{d-1}(q^h-1)\).
Distinct irreducible factors give distinct orbits. Since the
label is squarefree, parabolic induction of these characters is
the irreducible character indexed by the set of those orbits,
each with partition \((1)\); distinct labels give distinct characters.
The label function is precisely this induction product by construction.
For factor degrees \(d_j\), \(\sum_j d_j=r\), the inducing parabolic
has index \(\prod_{h=1}^r(q^h-1)/\prod_j\prod_{h=1}^{d_j}(q^h-1)\)
by the flag count. Multiplying by the inducing degree yields
\[
 \chi(1)=\frac{\prod_{h=1}^{r}(q^h-1)}{\prod_j(q^{d_j}-1)},
 \qquad \frac{\chi(1)^2}{|\mathrm{GL}_r(\mathbb F_q)|}
 \ge q^{-r}\prod_{h=1}^r(1-q^{-h})\ge w_q,
 \qquad w_q=(1-r/q)q^{-r}.
\]
Thus the positive coefficient at the identity of each character
idempotent is at least \(w_q\).

There are \(q^r\) monic polynomials of degree \(r\); at most
\((1+M_0)q^{r-1}\) have forbidden roots. A nonsquarefree
polynomial is \(p^2h\) with monic \(p,h\) and \(\deg h<r\);
for each possible degree of \(h\) at most \(q^{r-1}\) such
products occur. Hence the weighted count of allowed ones is
at least \(1-\eta\), where \(\eta=(1+M_0+2r)/q\).
There are \(\varphi(q^r-1)/r\) primitive polynomials: a primitive
element cannot lie in a proper subfield and has an orbit of
size \(r\) under Frobenius.
For an integer \(u\) with \(j\) distinct prime divisors, the
\(i\)-th smallest prime divisor is at least \(i+1\), giving
\(\varphi(u)/u\ge\prod_{i=1}^j i/(i+1)=1/(j+1)\);
also \(2^j\le u\). Consequently
\(u\le\varphi(u)(1+2\log u)\). With weight \(w_q\) for each monic polynomial, the count of nonprimitive polynomials is at most
\(\exp[-1/(2r(2r\log q+1))]\).
Taking the product of weighted counts over all \(f\) fields
and subtracting the tuples with no primitive polynomial gives
\[
 (1-e_i^{\rm pr})(1)\le1-(1-\eta)^f+
   \exp[-f/(2r(2r\log q+1))]\le f\eta+b^{-9}\le b^{-8}/2.
\]
The last two inequalities are the scale checks
\(0\le\eta\le1\), \(f\ge 2r(2r\log q+1)\,9\log b\) and
\(f\eta+b^{-9}\le b^{-8}/2\). The tensor product
\(e_0=e_0^{\rm pr}\otimes e_1^{\rm pr}\) acts as \(Q\).
This explains why both good-prime conditions are imposed.

For a complementary idempotent \(x=1-e_i^{\rm pr}\), the convolution
\(f(g)\mapsto\sum_{t\in G_i}x(t)f(gt)\) is an orthogonal projection
on the space of functions on \(G\) with operator-valued entries:
it is the group algebra action of a self-adjoint idempotent.
Apply this to \(E_g\) cut down to \(\Omega_i\). On \(R\)
its value is \(r_\psi(s)\sum_{t\in G_i}x(t)E_t\).
Using \eqref{B:nearBessel} and dividing by \(|R|\) gives
\[
 \left\|\sum_{t\in G_i}x(t)E_t\right\|_{\rm HS}^2\le b^{-12}d_X .
\]
For \(t\in G_i\) the other prime component is the identity.
A diagonal Bessel term is a rational element of \(U_{\rm Bes}\rtimes U(X_0)\)
whose conjugate by its class representative is integral and reduces
to the identity at that other prime. Here \(U_{\rm Bes}\) is the
unipotent radical of the Bessel subgroup. Its middle \(U(X_0)\)-block
inherits integrality and the identity reduction. By freeness at that
prime its rational part is the identity. Thus a diagonal term can
occur only if \(t\) is in the image of the integral unipotents. In
that case translating the identity block by a unipotent representative
multiplies each diagonal entry by its reduced character \(\psi_U(t)\).
Summing the \(d_X\) entries gives \(\operatorname{tr}r_\psi(t)=d_X\psi_U(t)\),
with \(\psi_U\) extended by zero elsewhere.
Centrality replaces \(\psi_U\) by its conjugacy average
\(\psi_U^\natural\) in the trace sum. The local two-frame mass
estimate implies
\begin{equation}\label{B:charmean}
 \sum_{g\in G}|\psi_U^\natural(g)|^2\le2 .
\end{equation}
This is input F5 below. The idempotent's self-adjointness and
idempotence give \(\sum_t|x(t)|^2=x(1)\). By Cauchy--Schwarz the
Bessel trace contribution has absolute value at most
\(d_X\sqrt{(b^{-8}/2)2}=b^{-4}d_X\).
The trace of the remainder has absolute value at most
\(\sqrt{d_X}\,b^{-6}\sqrt{d_X}\le b^{-4}d_X\).
Since \(\rho(x)\) is an orthogonal projection,
\[
 \|\rho(x){\Pi_{\rm sc}}\|_{\rm HS}^2
 =\operatorname{tr}({\Pi_{\rm sc}}^*\rho(x){\Pi_{\rm sc}})\le2b^{-4}d_X.
\]
The two good projectors commute. The squared norm missed by
their product is at most the sum missed by the factors
(orthogonal decomposition with respect to one factor).
Hence \(\|(1-Q){\Pi_{\rm sc}}\|_{\rm HS}^2\le4b^{-4}d_X\).
Subtract this from the inferior estimate following
\eqref{B:defect}, on any subspace \(Z\), to obtain \eqref{B:near}.

\paragraph{The global second moment.}
Set \(\beta(g)=\|q_g\|_{\rm HS}^2\). This nonnegative
function is of positive type: it is the Gram correlation of the
operators \(C_g\) from \S\ref{B:one}. In particular
\(\beta(g)\le\beta(1)\le2d_X\).

We give the finite group argument that converts the local
main term into the global bound. At a chosen place above \(p_0\),
a frame vector is a nonzero vector in the frame span, or in its
dual at the conjugate place. For a set \(T\) of such places,
let \(K_T\subset G\) fix one specified frame vector at each,
and let \(s_T=\sum_{g\in K_T}\beta(g)\). Define
\(D_{\rm fix}=b^{6k}\), \(C_T=D_{\rm fix}^{-|T|}I_*\), and \(x=b^{-4}\).
Call a reduction open if at some chosen place above \(p_0\) all
the relevant leading minors of the frame pairing matrix are
nonzero. The open-link estimate (F6) bounds the normalized
link moment on \(K_T\) intersected with the open reductions by
\((1+x)C_T\).
Since \(I_0\le C_T\) by \(b^{11kf}I_0\le I_*\), and
open reductions avoid scalar cosets, \eqref{B:rootbound} gives
\[
 \sum_{g\in K_T,\,g\text{ open}}\beta(g)
 \le (\sqrt{1+x}+b^{-15})^2 C_T\le(1+4x)C_T.
\]
For the fixer at all \(f\) places, the scalar-coset count is
at most \(M^f|R|\), each term at most \(2d_X\).
The scale bounds imply \(8M^f I_0\le C_T\):
use \(M\le b^{23}\), \(8\le b\), \(23f+1\le5kf\).
Off the scalar cosets, \(z\ge p_0^f\ge b^{10f}\);
the link bound is at most \(b^{-2kf}C_T\le C_T/8\),
and \(b^{-30}I_0\le C_T/8\).
Using \((u+v)^2\le2(u^2+v^2)\) in \eqref{B:rootbound},
the non-scalar contribution is at most \(C_T/2\), while the
scalar contribution is at most \(C_T/2\). Thus \(s_T\le C_T\)
when all places are fixed.

For other fixers a finite-field estimate is needed.
Let \(\Gamma\) be a finite group, let \(\beta:\Gamma\to\mathbb R_{\ge0}\)
be of positive type, let \(q\) be the residue field size and let \(\pi\)
be a representation of \(\Gamma\) in \(\mathrm{GL}_r(\mathbb F_q)\).
Suppose all stabilizers of nonzero columns of either of the
two injective frame maps \(A_0,A_1:\mathbb F_q^j\to\mathbb F_q^r\)
(with the action on the second frame contragredient) have
\(\beta\)-mass at most \(Y_0\).
Then
\begin{equation}
\label{B:minor}
 q\sum_{\det(A_1^t\pi(g)A_0)=0}\beta(g)
 \le 2\sum_{\Gamma}\beta(g)+q^{j+r/2}Y_0 .
\end{equation}
To verify this, take vectors \(u_g\) realizing positive type,
\(\langle u_a,u_b\rangle=\beta(a^{-1}b)\), and the delta and
Fourier orthonormal bases of \(\ell^2(\mathbb F_q^r)\).
The squared norm of
\(\sum_a\delta_{\pi(a)z}\otimes u_a\) is \(|\Gamma|\) times
the stabilizer mass of \(z\). The same identity in the Fourier
basis at \(\pi(a)^{-t}v\) holds for the dual stabilizer.
The cross inner product is
\(|\Gamma|q^{-r/2}\sum_g\psi(v^t\pi(g)z)\beta(g)\)
for a nontrivial additive character \(\psi\), with the Fourier
sign chosen accordingly. Cauchy--Schwarz bounds its real part,
so the weighted character sum has real part at most
\(q^{r/2}Y_0\) for \(v=A_1 y,z=A_0 x'\), both nonzero
parameter vectors. For \(Z_g=A_1^t\pi(g)A_0\), additive
orthogonality gives
\(\sum_{y\in\mathbb F_q^j}\mathbb E_{x'\in\mathbb F_q^j}
  \psi(y^t Z_g x')=|\ker Z_g^t|\), at least \(q\) if
\(Z_g\) is singular. Since \(\beta(g)\ge0\), multiply this pointwise lower bound by
\(\beta(g)\) and sum. The
axes \(y=0\) and \(x'=0\) cost at most twice the total mass;
the remaining pairs cost at most \(q^{j+r/2}Y_0\).
This proves \eqref{B:minor}. The weight in the application is a squared
Hilbert--Schmidt norm, so it satisfies the pointwise sign condition.

The sign condition is necessary in the general statement. Let \(S_3\) act on
\(\mathbb F_2^4\) by permuting the first three coordinates, and take its sign
character as \(\beta\). This one-dimensional character is of positive type.
Every vector and covector stabilizer contains a transposition, because two of
its first three coordinates agree. All their signed masses vanish. Set
\(j=2\), choose \(\operatorname{im} A_0=V=\langle e_1,e_2+e_4\rangle\),
and choose \(A_1\) so that \(\ker A_1^t=(12)V\).
The minor is singular exactly when \(gV\cap(12)V\ne0\).
Comparing the three nonzero vectors of each plane gives
\(g=1,(12),(123),(132)\). Their signs sum to two, whereas the total signed
mass is zero. With \(Y_0=0\), the claimed inequality without the sign
condition would read \(4\le0\).

Choose the set of places and their frame vectors maximizing
\(D_{\rm fix}^{|T|}s_T\); all choices are finite. At a place outside
\(T\), the mass in the current fixer of every additional
vector or covector stabilizer is at most \(s_T/D_{\rm fix}\), by
maximality. If \(T\) is full the preceding all-place bound
already suffices. Otherwise apply \eqref{B:minor} to each
leading minor at such a place; the relevant frame maps are
injective. The scale inequality \(p_0^{j+r/2}\le D_{\rm fix}\) holds for
the minor sizes \(1\le j\le\varrho+1\), and \(p_0\ge b^{10}\).
Each singular-minor mass is therefore at most \(3s_T/b^{10}\).
A non-open reduction has a vanishing relevant minor at this
place; \(\varrho+1\le b\) implies the mass of all non-open
reductions is at most \(xs_T\). Combining with the open estimate,
\[
 s_T\le(1+4x)C_T+xs_T\quad\Longrightarrow\quad
 s_T\le(1+9x)C_T
\]
because \(0\le x\le1/10\) and
\((1-x)(1+9x)-(1+4x)=x(4-9x)\ge0\).
The empty set of places is an admissible competitor, hence
\begin{equation}
\label{B:global}
 \sum_{g\in G}\|q_g\|_{\rm HS}^2
 \le D_{\rm fix}^{|T|}s_T\le(1+9b^{-4})I_* .
\end{equation}

\paragraph{The target dimension.}
By F7 the good forms on \(Y_0\) have the same dimension as those on
the terminal definite comparison space \(W_m\). The precise input is the
cuspidal trace comparison for the two inner forms with matching
Harish--Chandra parameters (F7 below): the rank is at least two,
the number of infinite places is even, and the required parameter
gaps hold. Appendix~\ref{sec:trace} proves the reduction to the trace comparison
and its local inputs for all permitted Hecke operators. Taking the identity
operator there and summing over the good types gives this dimension
equality. Thus \(d_T\) is at most the dimension of all
principal-level forms on the definite \(W_m\).
The level is free: a rational stabilizer element is integral at every place up to conjugacy and congruent to the identity at the chosen places above \(p_0\). The norm of its trace minus \(r\) is at most \((2r)^{2f}<p_0^f\). Divisibility by all the chosen primes forces trace \(r\), hence the identity by archimedean unitarity. Its index
in the maximal level is \(|G|\), so the class-mass formula (or
counting the free fibers of the double coset map with stabilizer
weights) gives \(|G|\mu\) classes.
Each contributes the dimension of the weight representation,
which is \(\delta\) by the Weyl degree formula~\cite[Lecture~24]{fultonharris} and the
parameter match in the admissible Bessel system:
for a definite rank-\(r\) representation of highest weight
\(\lambda\), the local degree is
\(\prod_{i<j}(\lambda_i-\lambda_j+j-i)/(j-i)\);
the discrete-series formal degree here has the same parameter
product at each place. Consequently
\(d_T\le |G|\delta\mu\).
Multiplying \eqref{B:global} by this upper bound cancels the
denominator \(\delta\mu\) in \(I_*\), proving \eqref{B:moment}.


\subsection{The cutoff error and the near-pair sums}
\label{B:cutoff}
The following calculation reduces the analytic cutoff error and
absolute convergence to a single geometric separation estimate, with
the fitted Bessel majorants of F0. Write \(h=\varrho+1\), \(i=r^3f\)
and \(d_0=2hrf\). For an archimedean tuple \(y\) put
\(u(y)=\prod_v\max(1,\ell(y_v)/b^{6h})\), with \(\ell(y_v)\) the
infimum of \(1+\|s y_v\|\) over the Bessel subgroup. Fix the class
\(c\), reduction \(g\), and point \(y\). Give a term leading to
class \(c'\) the weight \(\beta_{c'}(z)\), \(z=\gamma^{-1}y\).
If \(u(z)\le R\), \(R\ge1\), the frames in the support of this
cutoff lie in a ball of radius \(\rho=b^{37h}R\) in the
\(d_0\)-dimensional real space of their coordinates. The box property
and rigidity input F4 state that unequal rational frames have
coordinate distance at least \((N\rho)^{-1}\) in this ball.
A fixed frame has total weight at most one: two terms with the same
frame and reduction differ within a single class by a rational
unipotent stabilizer element, and its translates of the cutoff sum
to one. Translating one frame to the origin, and comparing disjoint
balls of radius \((2N\rho)^{-1}\) with the ball of radius
\(2\rho+(2N\rho)^{-1}\), gives weighted count
\[
 C(R)\le(1+4N\rho^2)^{d_0}
     \le (b^{6h^2}R^2)^{d_0}.
\]
Indeed \(N\le b^{60h}\), \(h\ge23\), and \(134h+1\le6h^2\).
Since \(4d_0<i\), layer integration gives
\[
 \sum_{c',\gamma}\beta_{c'}(z)u(z)^{-i}
   =i\int_1^\infty C(R)R^{-i-1}\,dR
   \le 2b^{12h^3rf}.
\]
For the local full and cut functions, the fitted majorants give
\(\|W_v^{\rm side}(x)\|u_v(x)^{3i}\le w_v(x)\le a_v\) and
\(a_v\int_{U_{\rm Bes}\backslash G_\infty(v)}w_v\le
 b^{400r^2}\dim V_v\). Each \(w_v\) is unipotent invariant.
On the locus where the two kernels differ,
\(u(y)u(z)\ge b\): some local cutoff acts at length
\(>b^{6h+2}/2\). For either kernel use \(2i\) of its spare powers at
each argument and \(i\) more at \(z\). The triangle inequality bounds
the difference by
\[
 2 b^{-2i}\prod_v w_v(y_v)a_v\ u(z)^{-i}.
\]
Multiply by \(\beta_c(y)\beta_{c'}(z)\), sum, and integrate. The
sum over \(c',\gamma\) costs at most \(2b^{12h^3rf}\), and integrating
\(\beta_c\prod_v w_v\) unfolds the unipotent variable with factor
\(\nu\). The row sum of norms of the differences is at most
\[
 4\nu b^{-2r^3 f+12h^3rf+400r^2 f}\prod_v\dim V_v
   \le \tfrac12\nu b^{-k^3f}.
\]
Indeed \(\prod_v\dim V_v\le b^{2kf}\), \(8\le b\), and
\[
 12h^3rf+1+400r^2f+2kf+k^3f\le 2r^3f.
\]
For the latter use \(12h^3+500\le k<r\): the terms other than
\(k^3f\) are at most \(402r^2f\le r^3f\). Without restricting to
the changed locus, the same majorant without the factor
\(2b^{-2i}\) bounds the full kernel and the cut kernel, proving
absolute summability of the integrals. Finally, row squared norms
are bounded by squared row sums, the number of classes is at most
\(I_0\), and \(|G|b^{40}\le4b^{2k^3f}\). The last estimate follows
from \(|G|\le M^{r^2 f}\), \(M\le b^{23}\) and
\(23r^2f+40\le2k^3f\). Summing rows and dividing by \(\nu^2\) gives
the error \(b^{-40}I_0\) used in \eqref{B:rootbound}.

\paragraph{Pairs with equal and unequal links.}\label{B:pairs}
A block squared is bounded by the corresponding sum over ordered
pairs of terms. For a pair in the same nonzero-discriminant link,
its first representative determines both the reduction and the
link, so summing all link squares incurs no extra multiplicity.
For the remaining terms expand the square after partition into the
nonzero-discriminant link blocks and one zero-discriminant block.
Keep the squares of the former and bound the remaining products by
the products of the term sizes, including the entire square of the
zero-discriminant block. The factor from two amplitude majorants
is \(\nu^2\) in either case.

For the arithmetic count let \(d\) be the pair discriminant,
\(j=k-\dim_E(\mathrm{common\ link})\), \(\tau=b^{kf/4}\), and
\(C_p>0\) the pair-count scale of the filling. The geometric
inputs bound the product of term sizes by
\[
 \frac{B}{\tau^{j+2}C_p d^k}
 \quad\text{if}\quad
 (d\eta)^2\le G_\gamma G_{\gamma'},\quad
 D'\ge\max(H_\gamma,H_{\gamma'}),\quad
 \tau^{j+6}D'/\eta^k\le B .
\]
Here \(G_\gamma\) is the level gauge and \(H_\gamma\) the cell-rank
moment, both for the indicated near representative. The
reciprocal-discriminant count is
\(\sum_{\text{pairs of corank }j}d^{-k}\le\tau^{j+2}C_p\).
There are at most \(h^3+1\le\tau\) coranks, so summation gives
\(\tau B\). For equal links one has \(j\le1,\ \eta=z,\ D'=I_*\):
the result is \(\tau^8 I_*/z^k=(b^{2f}/z)^k I_*\).
For the other pairs the required bound on the budget is
\(B=I_0/\tau^{10}\), giving \(\tau B\le b^{-kf}I_0\).
Section~\ref{app:pair-count} proves the weighted lattice count and reduces the
pair count to the one-term multiplicity over a prescribed lattice.
The gauge and cell-rank bounds on individual term sizes are also required.

\paragraph{The coefficient in the open-cell estimate.}
\label{B:openenergy}
For a fixed vector choice at the places in \(T\), take a finite
collection of nonzero open link blocks. The source's open-cell
parametrization assigns to each block a last entry \(t'\), a
representative \(\sigma\) in its floor group (the subgroup used in
the unipotent double-coset parametrization), and integral positions
with the prescribed fixer. It requires
\(|t'|_v^2>N^2\) at every infinite place. With the first
\(\varrho\) entries ranging over the finite entry set, the block
equals the sum of cut amplitudes of the double cosets represented
by \(w(p,t')\sigma\), with \(w(p,t')\) the Weyl representative
for this entry list. The class pair, reduction, and \(\sigma\)
determine the block, since \(\sigma\) determines the link. Thus one
may sum squares without a labeling multiplicity.

Write \(S_{\rm bl}\) for the sum of the block norm squares and
\(R_{k-1}>0\) for the filling's normalizing constant at predecessor
rank. The Siegel estimate after grouping the blocks by \(t'\) is
\[
 R_{k-1}^2 S_{\rm bl}\le(1+b^{-6})C_{\rm S} Z.
\]
Here \(C_{\rm S}\) is the scaled Siegel constant and \(Z\) sums,
over distinct last entries and pairs of heads, the real part of the
product of archimedean band-limit kernels times the local
last-entry density divided by the intermediate group's mass ratio.
The Gram identity at fixed \(t'\) supplies the kernel in this estimate.
For the intermediate group implementing the vector fixers, the
required sieve bound is
\[
 Z\le\left(\prod_{j=0}^{\varrho}C_{{\rm Pois},j}\right)
 \left[
  \left(\prod_{v}\frac{1+b^{-10}}{\deg_v}\right)
  \frac{1+b^{-6}}{b^{6k|T|}}+
  b^{2403(\varrho+1)^9 f}C_{\rm dec}
 \right],
\]
where \(C_{{\rm Pois},j}\) are the stage Poisson constants,
\(\deg_v\) are the local Bessel degrees, and \(C_{\rm dec}\) is
the filling decay constant. In particular the sieve contributes one
factor \(1+b^{-6}\) for the entire fixer set \(T\), including \(T=\varnothing\).
Under the measure and level identifications of \S\ref{spread:normalization}, the scalar calculation there bounds the displayed right side times
\((1+b^{-6})C_{\rm S}\) by
\(R_{k-1}^2(1+b^{-4})b^{-6k|T|}I_*\nu^2\).
Divide by \(R_{k-1}^2\) and take the supremum over finite collections
of nonnegative summands to get F6. Section~\ref{moment:open} gives the rational parametrization.
The fitted Gram/Siegel estimate, remaining local transform estimates and
identification of the measures are needed in \S\ref{moment:full}--\ref{moment:sieve}. The
\srcline{New/Spanning/MainTerm/EnergyEstimate}{631}{formal open-cell bound}
has leading coefficient \(1+b^{-4}\) when \(T=\varnothing\).
For \(T\ne\varnothing\) it has the additional factor \(b^{-6k|T|}\).
The leading coefficient matters for the dimension conclusion. For example,
let the trivial group act on an isometric inclusion of a \(d\)-dimensional
space into one of dimension \(2d\). Every translate has the same image.
The ratio of the second moment of the image projection to its
squared first moment, with trace normalized by \(2d\), equals
\((1/2)/(1/2)^2=2\). A moment estimate with this constant is compatible
with the translates filling only half the target.

\subsection{Arithmetic hypotheses for spreading}
\label{B:frontier}
For clarity, the following are hypotheses of the present
reconstruction until their own proofs are supplied. All filling
statements are for the particular admissible filling associated
to \(D\), with its prescribed local measures, ranks, levels and
Bessel data; they are not claims for arbitrary choices of cutoffs.
Appendix~\ref{sz:register} records which clauses are derived on the constructed scales.
\begin{enumerate}
\item[F0.] The construction of the admissible filling, including
the fitted Bessel systems, their finite-dimensional source
coordinates (nonempty index sets), formal degrees and scale
inequalities used above. In particular the unitary source rank is \(k\),
and the ratios involving \(I_0,I_*\) are part of the filling
admissibility, not consequences of linear algebra.
\item[F1.] The period-to-series identification at the fitted levels,
the normalized identity Bessel block, and the block change of
variables for a rational Bessel term with its unitary coordinate
matrix. These yield the equivariance and the trace calculation in
\S\ref{B:arithmetic}. Absolute convergence of the block integrals
follows from F0 and F4 by \S\ref{B:cutoff}.
\item[F2.] For any \(\Omega\) and positive integer \(z\) with
\(z\le\operatorname{meetNorm}(g)\) throughout \(\Omega\),
\(\mathcal E(\Omega)\le \nu^2(b^{2f}/z)^k I_*\).
This is a source-specific near-pair lattice estimate, not a
general second-moment bound.
\item[F3.] Expansion by links: if \(\mathcal E(\Omega)\le C\nu^2\),
the normalized cut moment of the remaining terms is
at most \(C+b^{-kf}I_0\). The missing proof is the cross-link
near-pair size sum, including zero-discriminant terms, and the
term classification on sets avoiding scalar cosets.
\item[F4.] For \(u(z)\le R\) the cutoff support lies in the
frame box of radius \(\rho=b^{37h}R\); two distinct rational frames
on that support are separated by \((N\rho)^{-1}\), and terms of a
fixed frame and reduction occur in a single unipotent orbit with
cutoff weights summing to at most one, as specified in
\S\ref{B:cutoff}. Section~\ref{moment:rigidity} gives the
integral coordinates, separation and orbit argument in the fitted model.
On the constructed scales that section also constructs the reduced cutoff and proves its partition identity for the levels and scales specified there. Thus F4 follows in the fitted model.
\item[F5.] The conjugacy-average bound \eqref{B:charmean} for the
specific reduced unipotent character. The local mass proposition
with congruence subgroup \(K(M)\) is
\[\int_K\left|\int_K\chi_M(yzy^{-1})\,dy\right|^2dz
\le2\operatorname{vol}(K)^2\operatorname{vol}(K(M)).\]
The function is constant on reduction cosets; integration
over these disjoint cosets followed by cancellation of their
positive volume proves \eqref{B:charmean}. Appendix~\ref{sz:shells} derives this integral bound from explicit
deep-vector frame and shell bounds. The primitive local frame
estimates stated there remain inputs.
\item[F6.] For open reductions in any fixer \(K_T\),
\(\mathcal E(K_T\cap\mathrm{open})\le
 (1+b^{-4})b^{-6k|T|}I_*\nu^2\).
Section~\ref{B:openenergy} reduces this sharp main-term estimate to the displayed Siegel and sieve inputs. The factor
\(1+b^{-4}\) cannot be replaced by an unspecified constant,
since the subsequent dimension argument depends on near-optimal
second moment.
\item[F7.] Use the two-inner-form trace input in \S\ref{sec:trace}
at the identity operator, the equality of the fitted Bessel
formal degree and the compact Weyl degree. These enter the bound on \(d_T\).
Lemma~\ref{lem:fitted-primary-character} supplies the finite character
normalization in the bad-type weight of \S\ref{B:arithmetic}.
\end{enumerate}


\subsection{Selection of one indexed family}
\label{B:selection}
By Theorem~\ref{A:main}, \(\mathcal S\subset\mathcal I\) spans \(S\).
Fix the translate set \(t\) obtained above. Linearity gives
\[
 F_t=\sum_{g\in t}P_g(S)
 =\operatorname{span}_{\mathbb C}\bigcup_{g\in t}P_g(\mathcal S).
\]
Successively choose an element of this union outside the span of the
elements already chosen until their span equals \(F_t\). Its finite
dimension ensures that the procedure ends. Retain a preimage and a
translate for each of the resulting basis vectors:
\begin{equation}
\label{B:witnesses}
 y_i=P_{g_i}\phi_i,\quad g_i\in t,\quad \phi_i\in\mathcal S,\quad
 1\le i\le n,\qquad (1-4/b^2)d_T\le n=\dim F_t\le d_T .
\end{equation}
Since \(1-4/b^2\ge 1/2\) and \(d_T>0\), this gives \(n>0\) and
\(d_T\le2n\). Each chosen preimage belongs to the integral lattice
and satisfies the shortness bound for every translate and conjugation.
Consequently the same witnesses meet both clauses of
Theorem~\ref{A:main} for the family \(y\).
